% The UA-PAR benchmark. Counts from paper/results/corpus_summary.json and the manifest; query rules from paper/guidelines/queries.md.
\section{The UA-PAR benchmark}\label{sec:uapar}

\subsection{Sources and corpus}
\uapar{} imitates a researcher's archive with public-source material in Ukrainian and Russian (licences in Section~\ref{sec:availability}). The corpus has 7\,070 files in six parts. \emph{Notes}: 2\,945 Wikipedia articles (1\,960 Ukrainian, 985 Russian) on 25 topics, many of them technical, among them antennas, radio, LTE, software-defined radio, signal processing, microcontrollers, cryptography, machine learning, physics, history and literature; each article is stored as a Markdown note with its source line. \emph{Documents}: 400 further articles (260 Ukrainian, 140 Russian) rendered to 133 PDF, 134 DOCX and 133 PPTX files, so the archive has office formats. \emph{Tables}: 150 data.gov.ua datasets, 100 kept as CSV and 50 converted to XLSX. \emph{Code}: 1\,775 files from 30 Ukrainian-topic GitHub repositories, indexed as a code collection. \emph{Pictures}: 1\,000 Crossmodal-3600 images. \emph{Speech}: 800 FLEURS recordings, 400 Ukrainian and 400 Russian. The articles were collected by topic search with a page-and-resume collector; the SDR topic did not have enough long articles, so the notes are 55 short of the planned 3\,000 rather than padded with weakly related articles.

Not every collected file is indexable: archive-index skips generated and vendored files in repositories, so the full EG2 index holds 6\,054 documents (759 from the code part) and 96\,908 chunks. The text-only E5 index holds the 4\,137 text documents of the same corpus and no pictures or speech; Section~\ref{sec:results} shows that restricting EG2 to the same 4\,137 documents changes its nDCG@10 by $-0.0015$.

\subsection{Queries}
There are 200 queries in five categories with fixed quotas: 40 \emph{known-item} queries describing a specific document without quoting its title; 60 \emph{topical} queries that several documents may answer; 40 \emph{technical} queries containing identifiers, model numbers or standards as people write them; 30 \emph{cross-lingual} queries in Russian for a Ukrainian target or the reverse; and 30 \emph{no-answer} queries, plausible questions about things the corpus does not contain, often with invented model numbers. 126 queries are in Ukrainian and 74 in Russian. The 170 answerable queries were each written for one target note, which is recorded but is not the only relevant document. A validator checks the quotas and language shares, that every target exists, that no topical or cross-lingual query shares a run of more than five consecutive tokens with its target (no copying), and that every no-answer query has zero lexical matches in the index. All 200 queries pass.

The queries were written by GPT-6-Luna through the Codex CLI under these rules. Because the query writer must not grade its own queries, GPT grades never enter the relevance labels; they are used only to compare grading behaviour, which does not separate self-preference from general leniency (Section~\ref{sec:results}). The no-answer queries have no lexical match by construction, so BM25 returns nothing for them; the comparison of systems on these queries is therefore a check of whether a dense system ever abstains, not a fair contest between lexical and dense retrieval.

\subsection{Relevance assessments}
The assessment pool is the union of the first ten documents returned by each of the ten systems: 5\,572 query--document pairs for 200 queries. Separate shuffled sheets hide the retrieval system identities. Each judge sees the query, document title and beginning, limited to 1\,200 characters; speech records include the source transcript. Grades are 0 (unrelated), 1 (partly relevant) and 2 (directly answers the query). These labels describe the visible snippet rather than an inspection of the complete document.

The initial primary judge is Gemini~3.8 Flash, invoked through \code{agy} in batches of 50. Gemma~4 12B through Ollama provides the main sensitivity assessment, and Gemma~4 E4B provides an ancillary open-weights assessment. GPT-6-Luna authored the queries; its grades are reported only for comparison and never enter the primary or sensitivity relevance labels. All initial sheets are preserved. Gemini and Gemma share a provider, so separate grading does not establish independence of training data or errors.

Initial Gemini and Gemma~4 12B disagree on 1\,091 pairs. For the reviewed variant, a separate Gemini pass reassesses those pairs with both initial grades visible under neutral A/B labels, balanced in presentation order. Agreed grades are retained. This is a within-model review, not independent third-judge adjudication. Retrieval results are compared across initial Gemini, initial Gemma~4 12B and reviewed Gemini. The choice of the primary judge followed inspection of agreement and target/no-answer diagnostics and preceded the reported aggregate effectiveness comparison of systems; it was not preregistered. No completed human relevance audit is claimed.

All 170 answerable query targets are notes, and 134 targets occur in the depth-10 pool. Consequently, target-grade diagnostics use a denominator of 134 and do not measure accuracy on the complete target set. The 30 designated no-answer queries are analysed separately.

\subsection{Multimodal subsets}
Text-to-image and text-to-speech retrieval are evaluated outside the pooled benchmark, using caption--image and transcript--recording pairings from the source datasets rather than separately collected relevance judgements. For pictures, the query is the first Ukrainian or Russian Crossmodal-3600 caption of an image, and the gold item is that image among all 1\,000. For speech, the query is the FLEURS transcript of a sentence and the gold items are the recordings of that sentence; FLEURS is n-way parallel, so a Russian transcript can query Ukrainian recordings of the same sentence and the reverse. Each condition is scored three ways: against the audio vector of a recording, against its transcript embedded as text, and against the better of the two. When the query is the transcript of the recording itself, the text and combined scores are an upper bound, not a retrieval result; the audio-only and cross-lingual conditions carry the claims.
